% Papiers de Sophie Germain — Français 9118, lot 5, vues Gallica 81 à 100.
%
% Pass: Opus 5 (claude-opus-5), 11 September 2026 — first pass, unchecked against
% the views by a human.
%
% What the twenty views turned out to hold. Three things, and not one letter,
% although 9118 is the correspondence volume:
%
%   * Views 81-85 (ff. 74-77, and p. (8) facing f. 78): an eight-page printed
%     pamphlet, an offprint headed « Extrait de la Bibliothèque française,
%     III.me année, N.o III », carrying Villoison's Latin birthday poem for
%     Lalande with the editors' French footnotes. The pamphlet is in her papers
%     because it is partly about her: footnote (4), on p. (6), names
%     « Mademoiselle Sophie Germain » and says she excels in mathematics, and
%     the verse it hangs on — « Cui jam Germanæ facies invisa peritæ » — puts
%     her in the poem. One verse on p. (7) carries an ink correction on this
%     copy: « atque » struck, « ipsa » written above.
%
%   * Views 85, 87, 89 (ff. 78, 80, 82): three receipts of the Institut, all
%     signed Cardot, for memoirs entered in the prize competitions. 21 September
%     1811 and 21 September 1813, both for the question on the vibrations of
%     elastic surfaces, with the Newton and Bacon epigraphs; 8 November 1821,
%     for the Académie's mathematics prize for 1822, with an epigraph from
%     Gauss. Each is dated, which almost nothing else in this corpus is. The
%     1821 slip carries a line asking « Monsieur Le chevalier de Karcher » to
%     pass the receipt to the author of the memoir.
%
%   * Views 91-99 (ff. 84-93): loose leaves in her hand. Two slips work the same
%     problem, headed « fig 39 » — the centre of gravity of a triangle by
%     balancing a parallelogram against two half-size triangles; f. 84 is a
%     short first go, ff. 92r-92v a fuller one that reaches x = a/3. Three
%     reading notes on number theory, each keyed to a page of a book she is
%     reading: p. 217 on p^4 + 4, p. 220 on p^4 + q^4 and the Fermat numbers,
%     and on f. 86 an objection to Gauss, Disquisitiones p. 272, on the
%     composition of binary quadratic forms — she says the product of two forms
%     always gives two forms, that Gauss seems not to have remarked it, and that
%     several things in his theory would then need rectifying. Three leaves in a
%     neater, rounder, more schoolroom hand on physics: Abeille holding cold to
%     be a real being, a reply to Abeille on the fire and light of « mon
%     opinion », and « Idées sur les Sens ». One leaf, f. 93, on Euler's § 45
%     and the nodal lines of a vibrating lamina, in which she opposes to him
%     what she has found herself.
%
%   Blank, and therefore absent below: views 86, 88, 90, 97, 100.
%
% Notation. Her second differential is set as she writes it, ddy/ds^2 with the
% straight d of Leibniz, not the round d; the first is dy/ds. Ternary forms are
% given in her own bracketing, (16, 3, 19). The printed pamphlet's page numbers
% are cited as « p. (5) » and its footnote calls as (1), (2), … The paragraph
% sign in the reference to Euler is set as §, which is what the leaf draws.
\documentclass[11pt,a4paper]{article}
% The preamble shared by every transcription of the mathematicians' manuscripts in Gallica.
%
% It rests on one idea: **uncertainty compounds, it does not resolve.** A
% transcription that smooths over a doubtful word has destroyed the only thing
% it was made for — a reader must be able to tell, at every point, what was on
% the page and what was supplied. The apparatus macros below are therefore the
% critical apparatus, not decoration.
%
% The same commands are recognised by `scripts/render.mjs`, which produces the
% site's reading view, and by `scripts/tei.mjs`. A macro added here must be
% added there too, or rendering fails loudly — which is the wanted behaviour.
%
% Comments here are English, like the rest of the repository. The documents
% this preamble sets are French, and that is deliberate: see the skills.

\ProvidesPackage{germain}[2026/09/15 Transcriptions of mathematicians' manuscripts in Gallica]

\usepackage{amsmath,amssymb,amsthm}
\usepackage{mathrsfs}
% Kept from the parent preamble so the renderer's diagram support stays
% available; these manuscripts draw no commutative diagrams, and nothing here uses it.
\usepackage{tikz-cd}
\usepackage[english,french]{babel}
\usepackage{fontspec}

% --- Greek ------------------------------------------------------------------
% The text face stays fontspec's default, Latin Modern, which has no Greek:
% XeTeX drops every glyph it lacks with a line in the log and nothing on the
% page, and the Greek words the manuscripts quote vanished from the PDFs. So
% the Greek and Greek Extended blocks get a XeTeX character class of their own,
% and entering or leaving a run of it switches the family to CMU Serif — the
% same Computer Modern design, with polytonic Greek — and back. Only the
% family changes: a Greek word in \textit or \textbf keeps its shape and series.
% The transitions to and from every other class carry the tokens that class
% has towards an ordinary letter, so French spacing before « ; » or inside
% « » still holds after a Greek word. No grouping, so no brace or \endgroup
% can come out unbalanced; maths is left alone, since XeTeX inserts nothing
% there. Kept identical in `exercices.sty`.
\makeatletter
\newfontfamily\germainGreekfont{cmun}[
  Extension=.otf, NFSSFamily=germaingreek,
  UprightFont=*rm, ItalicFont=*ti, SlantedFont=*sl,
  BoldFont=*bx, BoldItalicFont=*bi, BoldSlantedFont=*bl]
\newXeTeXintercharclass\germain@greek
\def\germain@greekrange#1#2{%
  \@tempcnta=#1\relax
  \loop
    \XeTeXcharclass\@tempcnta=\germain@greek
    \ifnum\@tempcnta<#2\advance\@tempcnta\@ne\repeat}
\germain@greekrange{"0370}{"03FF}
\germain@greekrange{"1F00}{"1FFF}
\def\germain@greekprev{\rmdefault}
\def\germain@greekin{%
  \edef\germain@greekprev{\f@family}\fontfamily{germaingreek}\selectfont}
\def\germain@greekout{\fontfamily{\germain@greekprev}\selectfont}
\def\germain@greekjoin{%
  \XeTeXinterchartokenstate=\@ne
  \@tempcnta=\z@
  \loop
    \ifnum\@tempcnta=\germain@greek\else
      \XeTeXinterchartoks\germain@greek\@tempcnta=\expandafter{%
        \expandafter\germain@greekout\the\XeTeXinterchartoks\z@\@tempcnta}%
      \XeTeXinterchartoks\@tempcnta\germain@greek=\expandafter{%
        \the\XeTeXinterchartoks\@tempcnta\z@\germain@greekin}%
    \fi
    \ifnum\@tempcnta<4095\advance\@tempcnta\@ne\repeat}
% babel-french sets its own classes, and saves and restores the interchar
% state, each time a language is selected: join again after it does.
\addto\extrasfrench{\germain@greekjoin}
\addto\extrasenglish{\germain@greekjoin}
\AtBeginDocument{\germain@greekjoin}
\makeatother

\usepackage{geometry}
\geometry{a4paper,margin=25mm,marginparwidth=28mm}
\usepackage{marginnote}
\usepackage{xcolor}
\usepackage[normalem]{ulem}
\setlength{\emergencystretch}{3em}
\usepackage{hyperref}
\hypersetup{colorlinks=true,linkcolor=black,urlcolor=black,pdfborder={0 0 0}}

% --- Batch metadata ---------------------------------------------------------
% Taken as they stand by the rendering script, which shows them at the head of
% the reading view. They fix what the file claims to cover. The macro names are
% the parent project's, so that the scripts stay shared; \folder here means the
% volume's slug — `fr-9115` — and \pages counts Gallica views.

\newcommand{\folder}[1]{\gdef\grothFolder{#1}}
\newcommand{\batch}[1]{\gdef\grothBatch{#1}}
\newcommand{\pages}[2]{\gdef\grothFirst{#1}\gdef\grothLast{#2}}
\newcommand{\foldertitle}[1]{\gdef\grothFoldertitle{#1}}

% The holder's own shelfmark, printed as they write it: « Français 9115 ».
\newcommand{\shelfmark}[1]{\gdef\grothShelfmark{#1}}

% The dating, copied verbatim from the holder's catalogue — never inferred
% here. « XIXe siècle » is the BnF's; a pass that dates a leaf from its content
% says so in a \note{}, as its own inference.
\newcommand{\dating}[1]{\gdef\grothDating{#1}}

% The status notice, printed in the title block and repeated as a diagonal
% watermark on every page of the PDF. The manuscripts are in the public
% domain; what this line declares is not a right but a quality — an
% unverified first pass by a machine. The skills write the line; a file
% without it gets no watermark, so leaving it out is visible in the output.
\newcommand{\watermark}[1]{\gdef\grothWatermark{#1}}

\gdef\grothFolder{?}\gdef\grothBatch{}\gdef\grothFirst{?}\gdef\grothLast{?}%
\gdef\grothFoldertitle{}\gdef\grothShelfmark{}\gdef\grothDating{}\gdef\grothWatermark{}

% --- The résumé -------------------------------------------------------------
\newenvironment{resume}
  {\small\setlength{\parskip}{0.45em}}
  {\par\medskip\hrule\medskip}

% --- Keywords ---------------------------------------------------------------
% In English, closing the modernised reading's résumé: the modern vocabulary
% under which to search for what the leaves construct. The manifest extracts
% them to tag the volume — this line is the single source of the tags.
\newcommand{\keywords}[1]{\par\smallskip\noindent
  {\footnotesize\textsc{Keywords} --- \textit{#1}}\par}

% --- The critical apparatus -------------------------------------------------

% The Gallica view number, in the margin. This is the anchor that lets a
% disagreement over a formula be settled by looking at *one* image rather than
% a volume of 750. The site uses it to turn the facsimile as the transcription
% is scrolled. A view is one scanned image: usually one side of a leaf.
\newcommand{\page}[1]{%
  \marginnote{\footnotesize\color{gray!70}\textsf{v.\,#1}}%
  \label{page:#1}\ignorespaces}

% The BnF's pencilled foliation, where the leaf carries it and a pass can read
% it: « 198r ». This is what the literature cites — « ff. 198r–208v » — and
% the one line that turns a citation into a view. Recorded, never inferred:
% a folio a pass cannot read is not written.
\newcommand{\folio}[1]{%
  \marginnote{\footnotesize\color{gray!50}\textsf{f.\,#1}}\ignorespaces}

% The modernised reading's coarser counterpart to \page{}: which views a
% section is drawn from.
\newcommand{\pagerange}[2]{%
  \marginnote{\footnotesize\color{gray!70}\textsf{v.\,#1--#2}}%
  \ignorespaces}

% Illegible: never guessed.
\newcommand{\ill}{\textcolor{red!60!black}{[\,\ldots\,]}}

% A doubtful reading: the word is given, and flagged as uncertain.
\newcommand{\uncertain}[1]{\uline{#1}}

% An editorial addition — a missing letter, an implied word.
\newcommand{\add}[1]{\textcolor{blue!50!black}{[#1]}}

% Struck out by the writer. What was crossed out is part of the document.
\newcommand{\struck}[1]{\sout{#1}}

% The transcriber's note, on the state of the page or the mathematics.
\newcommand{\note}[1]{\par\smallskip{\small\color{gray!60!black}\textit{[#1]}}\par\smallskip}

% A marginal note of hers, distinct from the body of the text.
\newcommand{\marginal}[1]{\par\smallskip\noindent\hspace{1em}%
  {\small\color{gray!60!black}#1}\par\smallskip}

% --- Title ------------------------------------------------------------------
% The title block is French, like the documents themselves.

\AddToHook{shipout/background}{%
  \ifx\grothWatermark\empty\else
    \begin{tikzpicture}[remember picture, overlay]
      \node[rotate=52, scale=3.4, align=center, text=gray!18,
            font=\sffamily\bfseries]
        at (current page.center) {\grothWatermark};
    \end{tikzpicture}%
  \fi}

\AtBeginDocument{%
  \begin{center}
    {\large\bfseries Manuscrits de mathématiciens (Gallica) ---
      \ifx\grothShelfmark\empty\grothFolder\else\grothShelfmark\fi}\\[2pt]
    {\itshape\grothFoldertitle}\\[4pt]
    {\small vues \grothFirst--\grothLast
      \ifx\grothBatch\empty\else\ (lot \grothBatch)\fi}\\[2pt]
    \ifx\grothDating\empty\else
      {\small\color{gray!60!black}Datation du catalogue : \grothDating}\\[2pt]
    \fi
    \ifx\grothWatermark\empty\else
      {\small\bfseries\color{red!45!gray}\grothWatermark}\\[2pt]
    \fi
    {\footnotesize\color{gray!60!black}Transcription établie d'après les images IIIF de Gallica.
     Source gallica.bnf.fr / Bibliothèque nationale de France.\\
     Les lectures incertaines sont soulignées, les illisibles notées [\,\ldots\,],
     les ajouts entre crochets ; \textsf{v.} numérote les vues de Gallica,
     \textsf{f.} la foliotation au crayon de la BnF.}
  \end{center}
  \vspace{1em}}

\endinput


\folder{fr-9118}
\shelfmark{Français 9118}
\batch{5}
\pages{81}{100}
\dating{1801-1900}
\watermark{Lecture automatique\\non vérifiée}
\foldertitle{Correspondance de Sophie GERMAIN avec les mathématiciens et savants Cauchy, Delambre, Fourier, Gauss, Le Gendre, d'Ansse de Villoison, etc.}

\begin{document}

\section{Pièce imprimée : vers de d'Ansse de Villoison pour la naissance de Lalande}

\page{81}\folio{74}
\note{feuillet imprimé, page de titre de la plaquette, p. (1) ; le verso, p. (2), est blanc. Le timbre « Bibliothèque impériale, MSS. » est apposé dans la marge.}

VERS

DE JEAN-BAPTISTE-GASPARD D'ANSSE DE VILLOISON,

MEMBRE DE L'INSTITUT NATIONAL DE FRANCE,

Pour le jour de la naissance du célèbre Astronome Jérôme de Lalande ( le 11 juillet ).

\page{82}\folio{75}
\note{p. (3).}

EXTRAIT DE LA BIBLIOTHÈQUE FRANÇAISE, III.me ANNÉE, N.o III.

LES auteurs de la BIBLIOTHÈQUE FRANÇAISE dérogent, pour cette fois-ci seulement, à la double loi qu'ils s'étoient imposée, de ne parler que des ouvrages français, et de n'insérer que des analyses d'ouvrages imprimés ; mais le désir de rendre un hommage public au savant Villoison, l'un des hommes le plus justement célèbres de l'Europe, ainsi qu'à l'illustre Lalande, le père de l'astronomie en France, suffit sans doute pour les justifier de s'être écartés du plan qu'ils se sont fait, et que le public a paru approuver. D'ailleurs, ils ont saisi avec plaisir cette occasion de rappeler à la France savante le nom de quelques-unes des personnes célè-

\page{83}\folio{76}
\note{la vue montre les pages (4) et (5) en regard.}

bres qui embellissent la cour du Duc régnant de Saxe Weymar.

VERS DE JEAN-BABTISTE-GASPARD D'ANSSE DE VILLOISON, Membre de l'Institut National de France, pour le jour de la naissance du célèbre Astronome Jérôme de Lalande ( le 11 juillet. )
\note{« BABTISTE » est bien ce qu'imprime la p. (4) ; la page de titre, p. (1), imprime « BAPTISTE ».}

\emph{Genethliacon Hieronymi Landii} ( DE LALANDE ), \emph{clarissimi astronomi.}

SIDERA concelebrant hodiernum sideris ortum :

Lætius insolito nunc vestit lumine cœlum

Landius exoriens, totumque amplectitur orbem.

Hoc nascente novus fastorum nascitur ordo ;

Inde notent annos et signent tempora docti.

Vestra est ista dies : Musarum plaudite alumni.

Pierioque choro et formosis dulce puellis

Si tellus nomen taceat, resonabit Olympus.

Irrita sacrilegæ requiescant murmura linguæ.

Niliacas quondam ad ripas, gens torrida solem

Ignivomum increpitans, voce adlatrabat inani,

Infelix rana, atque impar congressa, coaxat !

Gentem despiciens penitus penitusque jacentem,

Phœbus inexhaustæ fundebat flumina lucis,

Obscuram illustrans flammis ultricibus oram.

Sideribus cognatam animam formavit Apollo,

Ardoremque suum, et divinæ semina mentis

Indidit : Uraniæ patuit certissima proles.

Ardens flamma petit flammati culmina cœli.

Terram habeant reges : noster sibi servat Olympum.

Crasso ficta luto, atque hominum conspersa cruore,

Illi sordet humus cui ridet purior æther.

Impatiens volucri contendit ad æthera cursu,

Duxerat unde genus ; trahit hunc cœlestis origo :

Evehit, et celeres vigor igneus adjicit alas.

Surgentemque nepos sequitur, (1) conjuxque nepotis (2) ;

\note{notes du bas de la p. (5).}

(1) Français-de-Lalande, membre de l'Institut national de France, savant modeste, et digne neveu de l'illustre astronome de ce nom, a donné la Théorie de Mars, et un catalogue de cinquante mille étoiles, le plus énorme travail que jamais aucun astronome ait présenté. « C'est \emph{le tableau le plus fidèle et le plus complet qu'on pût désirer de l'état du ciel à la fin du dix-huitième siècle ; tableau qui ne fera que gagner de plus en plus en vieillissant, et que les siècles à venir citeront plus souvent, et avec plus d'éloges que les contemporains de l'auteur,} » dit l'immortel Delambre, dans son rapport lu à l'Institut de France.

(2) Madame Français-de-Lalande, dans la fleur de l'âge

\page{84}\folio{77}
\note{la vue montre les pages (6) et (7) en regard.}

Burchardusque (3) simul procedit passibus æquis.

Cui jam Germanæ (4) facies invisa peritæ,

Vidit, et invidit, cedens Ariadna coronam.

« Quæ nova stelligeris succedit sedibus hospes,

Erigone exclamat ? « Confidentissima nostras

» Tentat adire domos ! Divi prohibete volantem :

» Dum licet, Icariam superi frænate puellam ;

» Namque giganteos superabit fervida nisus.

\note{notes du bas de la p. (6) ; la première reprend la note (2) de la page précédente.}

et de la beauté, passe les nuits à observer les astres avec son mari, et destine ses enfans à la même carrière, où elle s'est couverte de gloire.

(3) M. Burckhardt, né dans le duché de Saxe-Gotha, demeure chez les savans de Lalandes, et s'est rendu justement célèbre par une foule d'observations astronomiques, \uncertain{et dernièrement encore par son beau travail sur la dixième} planète découverte par M. Olbers, médecin de Brême. On disputoit pour savoir si ce n'étoit pas une comète : M. Burckhardt a levé ce doute, et a calculé les dérangemens que Jupiter cause à cette planète, qui ne paroît que comme une étoile de huitième grandeur.
\note{la ligne soulignée est à demi rognée par le bord du cadre du microfilm ; la lecture s'appuie sur ce qui subsiste des lettres.}

(4) Mademoiselle Sophie Germain excelle dans les mathématiques. J'espère que sa modestie voudra bien me pardonner d'avoir trahi son secret. Qui pourroit résister à la tentation de lui rendre justice, ainsi qu'à mademoiselle Armande le Boulanger, digne sœur de madame de Beaumont, et pleine de goût et de connoissances dans la littérature grecque et latine ?

\note{p. (7).}

» Fœmina jam Placei regnum ambitiosa pererrat (5) !

» AEmula nunc aquilæ fertur Cytherea columba,

» AEthereosque haurit cupidis obtutibus ignes !

» Audax attonitis pulchrum caput inserit astris,

» Percurritque polum. Cœlo terraque marito

» Hæret juncta comes ; cunctis nunc devovet astris

» Quas Veneri noctes meliorem debet in usum,

» Immemor ; \struck{atque} ipsa viam natis quâ se quoque possint
\note{correction manuscrite portée à l'encre sur cet exemplaire imprimé : « atque » est biffé et « ipsa » écrit au dessus.}

» Tollere humo, patriumque sagax præmonstrat Olympum :

» Atque jubet similes stirpi succrescere ramos,

» Semper fronde novâ, et redivivo flore virentes,

» Quæ diffusa, ingens, mox cœlum conteget arbor.

» Fulminis hæredes quondam laudisque futuros,

» Sic avium regina suos educere fœtus

» Gestit, et implumes magnis jam destinat ausis.

» Tanta tenet tanti generis fiducia matrem !

» Huic tincta ingenio scintillant lumina : doctos

» \uncertain{Landius huic radios, pomum Paris ipse dedisset} ;
\note{ce vers aussi est à demi rogné par le bord du cadre.}

» Huic quoque diva favet Gotthæ (6) quæ præsidet arci,

» Cui Bereniceos Conon tribuisset honores.

\note{notes du bas de la p. (7).}

(5) Le sénateur Laplace, membre de l'Institut, que son génie sublime met au-dessus de tous les titres et de tous les éloges.

(6) S. A. S. madame la duchesse régnante de Saxe-Gotha, protége et cultive avec succès l'astronomie, et a fait un accueil distingué au patriarche des astronomes, et à sa docte

\page{85}\folio{78}
\note{la vue montre à gauche la dernière page imprimée, p. (8), qui achève la note (6), et à droite le récépissé de 1813, qui porte au crayon le folio 78 et ouvre la suite.}

nièce. Tout le monde sait que les princes des différentes branches de la maison de Saxe, sont les Médicis de l'Allemagne, et que le délicieux séjour de M. Wieland, de M. Goethe, de M. Herder, de M. Schiller, de M. Kotzebue, de M. Bœttiger, de M. Schwabe, etc., de mesdames les baronnes de Wolzogen, et Frédérique de Riedesel, de mademoiselle Amélie d'Imhoff, de madame Sophie Mereau, etc., etc., la ville de Weimar, est la Florence, ou plutôt l'Athènes de nos jours, comme Iena en est l'Alexandrie, graces aux soins éclairés et vivifians de LL. AA. SS. le duc régnant, et de mesdames les duchesses régnante et douairière de Saxe-Weimar.

\section{Récépissés de l'Institut}

\note{le premier des trois occupe la moitié droite de la vue 85, f. 78. Il est écrit sur un formulaire à en-tête calligraphié.}

Institut Impérial de France.

Paris, le 21 Septembre 1813

J'ai reçu un Mémoire pour concourir au prix proposé par la première classe de l'Institut sur cette question : \emph{Donner la théorie mathématique des vibrations des surfaces élastiques et la comparer à l'expérience}, ayant pour épigraphe ce passage de Bacon.

\emph{Sed longe maximum progressibus Scientiarum et novis pensis ac provinciis in eis denuo suscipiendis obstaculum deprehenditur in disperatione hominum, et suppositione impossibilis.}

que j'ai numéroté I et dont j'ai délivré le présent récépissé.

Cardot, employé à l'Institut.

\page{87}\folio{80}

Institut Impérial de France.

Paris, le 21 Septembre 1811

J'ai reçu un Mémoire pour concourir au prix proposé par la 1ère classe de l'Institut sur cette question : \emph{Donner la Théorie mathématique des Vibrations des Surfaces élastiques et de la comparer à l'expérience} ayant pour épigraphe : \emph{Effectuum naturalium ejusdem generis eædem sunt causæ.} Newton philo. nat. que j'ai numéroté I et dont j'ai délivré le présent récépissé.

Cardot, Employé au Secrétariat de l'Institut
\note{la question du prix est copiée ici « et de la comparer à l'expérience », et sur le récépissé de 1813 « et la comparer à l'expérience ».}

\page{89}\folio{82}

Institut Royal de France.

Paris, le 8 Novembre 1821.

J'ai reçu un Mémoire pour concourir au prix de Mathématiques proposé par l'Académie Royale des Sciences pour l'année 1822, \emph{sur le meilleur ouvrage, ou Mémoire de Mathématiques pures ou appliquées, qui aura paru, ou qui aura été communiqué à l'Académie, dans l'espace des deux années qui sont accordées aux concurrents}, ayant pour épigraphe :

« \emph{Illecebris harum quaestionum ita sui implicatus ut eas deserere non potuerim.} Gauss.

que j'ai numéroté I, et dont j'ai délivré le présent récépissé.

Cardot, Chef du Bureau du Secrétariat de l'Institut Royal

Monsieur Le chevalier de Karcher est prié de vouloir bien faire passer ce récépissé à l'auteur du Mémoire ci-dessus.

\section{Feuillets détachés : centre de gravité, notes de lecture, physique}

\page{91}\folio{84}
\note{la vue montre deux billets superposés, f. 84 et f. 85, sans rapport l'un avec l'autre. Le premier est une première rédaction du problème que reprend le f. 92.}

fig 39 Le triangle $ABC$ = parallélogramme $AFDE$ + 2 triangle $EDC$ = 2 parallélogramme $AEDF$,

Le centre de gravité du parallélogramme est, comme on sait d'ailleurs en $c'$ milieu de la diagonale $AD$, on a évidemment $EE' = FF' = \frac{1}{2} AD$

Les centres de gravité des triangles $EDC$, $EDB$, sont évidemment sur les lignes $EE'$ $EE'$ et dans une ligne parallèle à la base $BC$.
\note{la seconde ligne est écrite $EE'$ ; le sens demande $FF'$.}

Soit $EE' = FF' = a$ et par conséquent $AD = 2a$

\folio{85}
\note{second billet de la même vue, d'une autre affaire : une note de lecture, renvoyant à la p. 217 d'un ouvrage qu'elle ne nomme pas.}

p. 217

Aucun nombre de la forme $p^4 + 4$ excepté 5 n'est un nombre premier

car $p^4 + 4 = (p^2 - 2)^2 + 4p^2$ et par conséquent ces nombres sont, de plusieurs manières la somm\add{e} de deux quarrés
\note{la fin du mot passe hors du cadre.}

Pour $5 = 1 + 4$ $(p^2 - 2)^2 = (1 - 2)^2 = 1$ et les deux membres sont identiques

\page{92}\folio{86}

Gauss pag. 272 dit que le produit des deux formes $(16, 3, 19)$ $(8, 1, 37)$ est $(8, -1, 37)$

Ce produit est aussi $(13, -2, 23)$ \struck{Comme} on peut s'en assurer

En général il y a toujours deux formes qui résultent du produit de deux formes données, mais dans les cas particuliers les deux formes se réduisent quelquefois à une.

Il semble que l'auteur n'a pas fait cette remarque, et alors il y auroit plusieurs choses à rectifier dans la théorie qu'il donne du produit de plusieurs formes quadratiques binaires

\page{93}\folio{87}
\note{note de lecture sur la p. 220 du même ouvrage que le f. 85.}

p. 220

En général $p^4 + q^4 = (p^2 - q^2)^2 + 2(pq)^2 = (p^2 + q^2)^2 - 2(pq)^2$

Ainsi lorsque $p^4 + q^4$ est un nombre premier on peut être assuré, que ce nombre n'est que de ces seules manières la somme d'un quarré, \ill{} $+$ ou $-$ un double quarré. Les nombres $2^{2^i} + 1$ ne sont qu'un cas particulier de la forme $p^4 + q^4$
\note{deux ou trois mots biffés après « quarré, » sont illisibles ; l'exposant est écrit à deux étages, un $i$ au dessus du second 2.}

\page{94}\folio{88}
\note{d'ici au f. 90, trois feuillets d'une écriture plus ronde et plus posée que celle des billets de calcul, sur des sujets de physique.}

Monsieur Abeille pense que le froid est un être réel, il oppose à l'opinion de Physiciens qui ne le considère que comme l'absence du feu, le raisonnement suivant : l'eau qui se congèle augmente de volume, l'absence d'un être corporel, ne peut ajouter à la masse dont il se sépare, puisque le néant ne produit aucun effet positif, donc un être réel augmente le volume de l'eau.

La force avec laquelle l'eau augmente de volume en se congelant est bien prouvée par l'expérience suivante.

Que l'on prenne un tube de fer que l'on en ferme une extrémité, au moyen d'une virole de même Métal, qu'on \uncertain{emplie} ce tube d'eau, et qu'on l'expose au grand froid.

L'eau fera un tel effort qu'elle fendra le tube dans toute sa longueur, et l'on observera que les \uncertain{glaçons} déborderont cette fente.
\note{le feuillet paroît écrire « graçons » ; le sens ne laisse pas d'autre lecture.}

\page{95}\folio{89}

L'objection de Monsieur Abeille contre mon opinion sur le feu et la lumière est que l'on ne peut décomposer le feu, tandis que nous décomposons la lumière tant que nous le voulons.

Voici ma réponse je conviens que le feu n'a pas été décomposé jusqu'à présent, je crois même qu'il ne le sera jamais, mais je ne conclus pas de là que ce soit un être simple et je pense au contraire que s'il se refuse à devenir le sujet de nos expériences sous ce raport nous devons nous en prendre plutôt à l'impossibilité de l'obtenir seul, et de pouvoir le retenir ( puisqu'il pénètre tous les corps ) et au défaut de matière propre à opérer sur lui, qu'à la simplicité de son essence.

\page{96}\folio{90}

Idées sur les Sens

Je pense qu'ils peuvent être tous rapportés à celui du toucher et qu'ils ne diffèrent que par la Disposition à ressentir l'attouchement de corps de nature différentes

le sens du toucher proprement dit est le moins délicat de tous puisqu'il n'est guère affecté que par les corps les moins déliés

tandis que l'œil est sensible à l'attouchement du subtile élément de la lumière

que l'ouï reçoit l'impression des parties de l'air mises en vibrations par le corps sonore ou retentissant

que l'odorat distingue les parties déliées qui s'évaporent des corps et que le goût qui a une si grande affinité avec l'odorat connoit presque la forme des particules que le touchent par l'impressions qu'elles lui font ressentir ;
\note{la syntaxe des deux derniers membres est celle du feuillet.}

\page{98}\folio{92}
\note{reprise, plus complète, du problème du f. 84 ; la suite est au verso, que montre la vue 99.}

fig 39 Les côtés du triangle $ABC$ sont \struck{deux} égaux fois ceux des triangles $BFD$, $CED$ multipliés par 2, le centre de gravité du triangle $ABC$ est sur $AD$ à une distance $2x$ de la base $BC$, de même aussi les centres de gravité des triangles $BFD$, $CED$ sont sur $FF'$ $EE'$, à une distance $x$ des bases $BD$ $CD$, il faut déterminer la grandeur de $x$.
\note{« égaux » est écrit au dessus de « deux » biffé, mais « fois » n'a pas été biffé.}

Pour y parvenir nous observerons que la somme des forces de gravité des deux petits triangles est égale à celle du parallélogramme $AFDE$ et doit être placée sur la ligne $AD$ parallèle à $EE'$, $FF'$, dans le point où la parallèle à la base qui passe par les centres de gravité des petits triangles coupe $AD$.

\page{99}
\note{la vue montre à gauche le verso du billet de la vue 98, qui poursuit le calcul, et à droite un feuillet sans rapport, coté 93.}

Cette distance est $= x$, puisque les côtés des petits triangles sont moitié moins grands que ceux du triangle donné.

Si pour abréger on nomme $X$ le point où se réunissent les efforts des forces de gravité des petits triangles et $c'$ le centre de gravité du parallélogramme on aura $XC'$ sur la ligne $AD$ et par conséquent $=$ à la moitié $C'D$ on aura, à cause de l'égalité entre le poids du parallélogramme et la somme de ceux des petits triangles la moitié de $C'X$ pour le lieu du centre de gravité du triangle \ill{}

$C'X = \frac{1}{2}(a - x)$ on a donc $x + \frac{1}{2}(a - x)$
\note{la ligne se perd hors du cadre à droite ; dans « la moitié de $C'X$ » la seconde lettre est reprise sur une autre et sa lecture est douteuse.}

\marginal{$4x = a + x$ ; $x = \frac{a}{3}$ ; c. q. f. d. pour la distance de ce centre à la base et $3x = a$}
\note{la colonne marginale de gauche porte encore deux ou trois lignes de calcul, à demi biffées et illisibles.}

\folio{93}
\note{feuillet de droite : elle discute Euler sur les portions d'une lame vibrante de part et d'autre d'une ligne nodale.}

Euler établit \uncertain{suivant} ( § 45 de la \uncertain{P.} ) que les angles que l'une et l'autre portions font avec l'axe en $L$ doivent nécessairement être égaux et qu'aussi au point $L$ les rayons osculateurs de l'une et l'autre portions ne peuvent être différens, mais pour le cas où les extrémités sont libres ou fixés ces suppositions ne sont point valables pour des valeurs de $a$ différentes de $\frac{1}{2}$ puisque j'ai trouvé qu'alors les courbures prises à droite et à gauche du point $L$ étoient différentes entr'elles ce qui résulte de ce qu'aucune ligne nodale située ailleurs qu'au milieu de la lame n'est à distance égales des deux plus voisines

il semble aussi que $\frac{dy}{ds}$ $\frac{ddy}{ds^2}$ c'est à dire les tangentes et les rayons osculateurs doivent être de signes différens pour les deux portions
\note{le renvoi entre parenthèses est d'une lecture douteuse : le mot abrégé avant le paragraphe et la lettre après « de la » ne se laissent pas assurer.}

\end{document}
